% !TEX root = ../thesis.tex

\chapter{Experiments and Results} \label{chp:exp}

\section{Experimental Setup}

All experiments were conducted in NVIDIA Isaac Gym simulation with 2048
parallel environments running on a single NVIDIA RTX 4090 GPU.
The B2+Z1 URDF was loaded with the complete kinematic and collision meshes
for both the quadruped body and the Z1 arm.
Training was performed in three sequential stages as described in
Chapter~\ref{chp:method}: locomotion pre-training (50k iterations),
arm pre-training (30k iterations), and hybrid fine-tuning (30k iterations).
All rewards were logged via Weights \& Biases throughout training.

\section{Stage 1: Locomotion Results}

After 50,000 PPO iterations, the B2 quadruped (with arm fixed in rest
configuration) produced stable trotting gait on flat terrain.
Key observations:
\begin{itemize}
  \item The robot maintained an upright posture throughout 20-second rollouts
        without falling (terminal body-height condition never triggered).
  \item Forward velocity tracking error converged to within
        $\pm 0.1$\,m/s of the commanded value.
  \item The gait transitioned smoothly between different commanded
        frequencies ($2$--$4$\,Hz) and stance widths, confirming that
        the distributional command curriculum was active.
  \item Prior to the PD gain correction ($K_p=35$, $K_d=1$), the robot
        collapsed within the first few training iterations due to
        insufficient joint torque authority.  After the correction
        ($K_p=200$, $K_d=20$), stable locomotion emerged within the
        first 10k iterations.
\end{itemize}

\section{Stage 2: Arm Manipulation Results}

After 30,000 PPO iterations on the stand-still grasping task, the arm
policy learned to track randomly sampled end-effector pose targets while
the robot body remained stationary.
Key observations:
\begin{itemize}
  \item The arm successfully reached targets sampled from the full command
        range: reach length $l \in [0.3, 0.77]$\,m, polar angle
        $p \in [-0.45\pi, 0.45\pi]$\,rad, azimuthal angle
        $y \in [-\pi/2, \pi/2]$\,rad.
  \item End-effector orientation tracking converged; the weighted reward
        ($w_{\text{lpy}}=3$, $w_{\text{rpy}}=1$) was effective in
        prioritising position before orientation accuracy.
  \item The body orientation communication channel was active: the arm
        policy output non-zero pitch/roll commands to the dog policy when
        reaching targets near the workspace boundary.
\end{itemize}

\section{Stage 3: Hybrid Walk-While-Grasp Results}

The hybrid fine-tuning stage combined both pre-trained policies, enabling
the B2+Z1 system to walk while simultaneously tracking arm end-effector
targets.
Key observations:
\begin{itemize}
  \item The robot maintained stable locomotion at commanded velocities
        while the arm tracked moving end-effector targets.
  \item The arm policy successfully requested body lean adjustments via
        the orientation command channel to extend reachability.
  \item No sustained policy interference (locomotion destabilising the
        arm or vice versa) was observed in evaluation rollouts.
\end{itemize}

\section{Discussion}

The primary finding is that the RoboDuet cooperative architecture
transfers successfully to the heavier B2+Z1 platform once the PD gains
are properly calibrated.
The three-stage curriculum was essential: attempting to train all 18
joints simultaneously from scratch led to instability in preliminary
experiments.
Hot-starting Stage~3 from separately pre-trained checkpoints resolved
this issue.

\textbf{Limitations.}
The current work is simulation-only; sim-to-real transfer has not been
validated and would require additional domain randomisation tuning and
actuator identification for the physical B2 hardware.
The end-effector target is a fixed synthetic command; a full grasping
demonstration requires integration with a perception module to localise
real objects.

\chapter{Conclusion} \label{chp:conclusion}

This project reproduced the RoboDuet cooperative loco-manipulation
framework and adapted it to the Unitree B2+Z1 platform through systematic
identification and correction of platform-specific engineering mismatches.
The central contribution is a validated three-stage PPO training pipeline
--- locomotion pre-training, arm pre-training, and joint hybrid fine-tuning
--- that produces a B2+Z1 system capable of walking while tracking 6-DoF
end-effector pose targets in Isaac Gym simulation.

The cooperative dual-policy architecture proved effective on the heavier B2
platform: the body leans dynamically in response to arm-policy requests
without explicit user specification, replicating the key behavioural
property of the original Go1-based RoboDuet~\cite{pan2024roboduet}.

Future work directions include:
\begin{enumerate}
  \item \textbf{Sim-to-real transfer:} deploying the trained policy on
        the physical B2+Z1 hardware~\cite{lee2020learning,kumar2021rma}.
  \item \textbf{Object-centric grasping:} integrating a perception module
        to generate real end-effector targets from detected objects.
  \item \textbf{Rough terrain:} extending locomotion training to the
        trimesh terrain curriculum already in the framework.
  \item \textbf{Whole-body force control:} exploring compliant contact
        strategies for reactive manipulation.
\end{enumerate}
